\section{Forschungsmethodik}
### Forschungsfrage:  
IT-Governance als Einflussfaktor auf die Verwendung von Multifaktor-Authentifizierung in beruflichen und privaten Kontexten 

Hypothesen: 

H0: Die berufliche Verwendung von MFA hat keine Auswirkungen auf die private Verwendung. 

H1: Eine etablierte IT-Governance im beruflichen Kontext wirkt sich positiv auf die Nutzung und Akzeptanz von Multi-Faktor-Authentifizierung aus. 

H2: Eine hohe Benutzerfreundlichkeit von MFA-Systemen erhöht die Akzeptanz sowohl im beruflichen als auch im privaten Kontext. 

H3 (bzgl. Spill-over-Effekt): Positive Erfahrungen mit MFA im beruflichen Kontext erhöhen die Wahrscheinlichkeit der privaten Nutzung von MFA 

 

Methoden sind einerseits Deskriptiv, durch die Darstellung der Stichprobenverteilung und Nutzungsmuster, aber auch eine Inferenzstatistik zur Überprüfung der Zusammenhänge (Hypothesen) mittels Korrelationsanalyse (Pearson/Spearman). 
Umfrage, um zu quantifizieren, wie stark die berufliche Nutzung von MFA mit der privaten Nutzung korreliert, und ob Usability und Governance als Einflussfaktoren erkennbar sind.
 

#Entwurf Methodik 02/2026 

 ##1. Forschungsdesign und methodologische Einordnung 
Grundlegende Ausrichtung und Festlegung auf ein quantitatives Forschungsdesign. Für die Beantwortung der Forschungsfrage ist es entscheidend, die Stärke und Richtung von Zusammenhängen zwischen beruflicher MFA-Nutzung, Usability, Governance und privater Nutzung zu quantifizieren. Qualitative Methoden wie Interviews könnten zwar tiefere Einblicke in individuelle Erfahrungen bieten, wären jedoch weniger geeignet, um generalisierbare Aussagen über die breite Masse der Erwerbstätigen zu treffen. Daher ermöglicht die quantitative Online-Umfrage die Erhebung von Daten von einer größeren Stichprobe, was wiederum die statistische Analyse und die Ableitung von Kausalitäts-Tendenzen erleichtert. Zudem ist es wichtig, die Grundgesamtheit klar zu definieren (erwerbstätige Personen mit IT-Nutzung), um die Relevanz und Anwendbarkeit der Ergebnisse sicherzustellen.
 (erklären, warum Zahlen hier besser helfen als qualitative Interviews, z.B. um allgemeingültige Aussagen über Erwerbstätige zu treffen --> Rinke will immer wissen, was die Grundgesammtheit ist). 
Das Vorgehen ist deduktiv. Ableitung der Hypothesen aus bestehenden Theorien (Usability, Governance, Spill-over --> siehe Seminar) und deren empirische Überprüfung.



Wissenschaftstheoretische Verortung: Einordnung in den Kritischen Rationalismus (nach Karl Popper): Wir stellen Hypothesen auf und versuchen, die Nullhypothese (H0) zu widerlegen. 

### Forschungsziel und Fragestellung:  

Kurze Wiederholung der zentralen Frage: Wie beeinflusst der berufliche Umgang mit MFA das private Sicherheitsverhalten? 

Ziel: Identifikation von Kausalitäts-Tendenzen und Korrelationen. 

Abgrenzung des Untersuchungsraums: 

Fokus auf die Schnittstelle zwischen Arbeitswelt und Privatleben (Work-Life-Interface/Spillover). --> Methode hier wichtig, um den Umfang vom Spill-over zu quantifizieren 

Warum MFA (muss das hier wieder sein?)? (Aktualität durch Cyber-Bedrohungen und gesetzliche Vorgaben wie NIS-2). 

Begründung der Methodenwahl: Kombination aus systematischer Literaturanalyse (Theorie-Basis) und Online-Umfrage (Primärdatenerhebung). Vorteile der Online-Erhebung: Anonymität (wichtig bei Sicherheitsfragen), Erreichbarkeit einer großen Stichprobe, effiziente digitale Weiterverarbeitung in Orange/Python. 

## 2. Literaturanalyse (muss die sein?) 

Ziel der Analyse: 

Systematische Erfassung des aktuellen Forschungsstandes zu den drei Säulen: Usability, IT-Governance und Spill-over-Theorie. 

Identifikation von Forschungslücken (z. B. wenig Literatur zum Transfer von MFA-Verhalten vom Job ins Private). 

Methodik der Literaturrecherche: 

Suchstrategie: Festlegung von Schlüsselwörtern (Keywords) wie „Multi-Factor Authentication“, „User Acceptance“, „Spill-over Effect IT“, „IT-Compliance behavior“. 

Datenbanken: Nutzung einschlägiger wissenschaftlicher Quellen (z. B. Google Scholar, ResearchGate, ACM Digital Library, IEEE Xplore sowie die Bestände der Universitätsbibliothek). 

Einschluss- und Ausschlusskriterien: Warum wurden bestimmte Quellen gewählt? (z. B. Fokus auf Veröffentlichungen der letzten 10 Jahre, Peer-Reviewed-Artikel). 

Thematische Schwerpunkte der Aufarbeitung: 

Säule 1: Usability von MFA: 

Analyse bestehender Studien zur Nutzerfreundlichkeit (z. B. Zeitaufwand, kognitive Last bei verschiedenen MFA-Methoden wie SMS-Code vs. Authenticator-App). 

Bezug zum Technology Acceptance Model (TAM) oder UTAUT. 

Säule 2: IT-Governance im Unternehmen: 

Wie definieren Unternehmen Sicherheitsrichtlinien? 

Zusammenhang zwischen strengen Vorgaben (Compliance) und dem tatsächlichen Nutzerverhalten. 

Säule 3: Spill-over-Theorie: 

Theoretische Herleitung: Wie werden Gewohnheiten von einer Domäne (Professionell) in eine andere (Privat) übertragen? 

Bisherige Erkenntnisse zu „Security Spill-over“. 

Synthese und Hypothesenableitung: 

Zusammenführung der Erkenntnisse: Warum legen die gefundenen Quellen nahe, dass H1, H2 und H3 zutreffen könnten? 

Überleitung zum empirischen Teil: Die Literatur liefert die Theorie, deine Umfrage prüft die Praxis. 

3. Datenerhebung: Onlinefragebogen 

3. Datenerhebung: Quantitative Befragung 

Wahl der Erhebungsmethode: 

Begründung für die standardisierte Online-Befragung (Effizienz, Ortsunabhängigkeit, Anonymität zur Reduktion sozialer Erwünschtheit bei Sicherheitsthemen). 

Einsatz eines Umfrage-Tools (Microsoft Forms, weil ich damit arbeite). 

Fragebogenkonstruktion nach Porst (2014): 

Vorgehensweise: Anwendung der 10 Gebote der Frageformulierung nach Porst. 

Vermeidung von Fehlern: Erläuterung, wie du „doppelte Verneinungen“, „Suggestivfragen“ und „Überforderung“ vermieden hast. 

Fragetypen: Einsatz von geschlossenen Fragen und vor allem Likert-Skalen (z. B. 5-stufig von „stimme gar nicht zu“ bis „stimme voll zu“), um Einstellungen messbar zu machen. 

Operationalisierung (Messbarmachung der Variablen): 

Abhängige Variable: Die private Nutzung von MFA (Häufigkeit, Art der genutzten Dienste). 

Unabhängige Variablen: 

Usability: Wahrgenommene Einfachheit der Bedienung (basierend auf etablierten Skalen wie der System Usability Scale). 

IT-Governance: Wahrgenommener Druck oder Unterstützung durch betriebliche Richtlinien. 

Spill-over: Die psychologische Übertragung der Gewohnheit. 

Pre-Test (geht sich das noch aus?): 

Beschreibung eines Probelaufs (Pre-Test) mit einer kleinen Gruppe (z. B. 3-5 Personen). 

Ziel: Prüfung der Verständlichkeit der Fragen und der technischen Funktionalität. 

Dokumentation von Anpassungen nach dem Feedback (das füllt wertvolle Zeilen!). 

Stichprobenbeschreibung und Rekrutierung: 

Grundgesamtheit: Erwerbstätige Menschen in Deutschland/DACH-Raum mit IT-Nutzung. 

Stichprobenziehung: Gelegenheitsstichprobe (Convenience Sampling) über berufliche Netzwerke (LinkedIn/Xing) und private Kanäle. 

Datenschutz: Hinweis auf die Einhaltung der DSGVO, die Freiwilligkeit der Teilnahme und die anonyme Speicherung der Daten. 

Aufbau des Fragebogens (Dramaturgie): 

Einleitung (Zweck der Studie, Datenschutz). 

Demografie (Alter, Branche, IT-Affinität). 

Beruflicher Kontext (Governance-Vorgaben, MFA als Pflicht?) 

Private Nutzung (Status Quo) 

Einstellung/Usability-Bewertung 

Kommentare 

4. Hypothesen 

5. Datenanalysemethoden 

Datenanalyse und Auswertung mit OrangeWahl der Analysesoftware:Begründung für Orange Data Mining (Open-Source, visuelle Programmierung, Reproduzierbarkeit).Vorteil: Die Kombination aus explorativer Datenanalyse (Sichtprüfung) und statistischer Testung in einer Workflow-Umgebung.Datenimport und Vorverarbeitung (Data Preprocessing):File Widget: Import der bereinigten Daten aus der Online-Umfrage (CSV/Excel).Data Cleaning: Behandlung von fehlenden Werten (Imputation) oder Ausschluss von unvollständigen Datensätzen.Feature Selection / Edit Domain: Definition der Variablen-Typen (z. B. Ordinal für Likert-Skalen, Nominell für Geschlecht/Branche).Filterung: Ausschluss von "Ausreißern" (z. B. extrem kurze Bearbeitungszeit des Fragebogens).Explorative und Deskriptive Statistik:Distributions Widget: Visualisierung der Stichprobenverteilung (Alter, MFA-Nutzungshäufigkeit) mittels Histogrammen.Box Plots: Vergleich der Akzeptanzwerte (Usability) zwischen verschiedenen Nutzergruppen (z. B. IT-affine vs. weniger IT-affine Personen).Inferenzstatistische Analyse (Hypothesenprüfung):Correlation Widget: Berechnung der Korrelationen nach Spearman (da Likert-Skalen meist als ordinal skaliert betrachtet werden) zur Prüfung von $H_2$ (Usability) und $H_3$ (Spill-over).Scatter Plot: Visuelle Überprüfung der Zusammenhänge zwischen beruflicher Vorgabe (Governance) und privater Nutzung.Testung von Gruppenunterschieden: Einsatz des t-Test oder ANOVA Widgets, um zu prüfen, ob sich die private Nutzung bei Personen mit strenger beruflicher Governance signifikant von denen ohne unterscheidet ($H_1$).Dokumentation des Workflows:Erklärung, dass der gesamte Analyseprozess als Orange-Workflow (.ows) gespeichert wurde, was die Transparenz und Validität der Forschung erhöht (Replikationsfähigkeit). 

 

 

 